% =========================================================================
% Capítulo 2: Cinética de la partícula
% Tema 01: Dinámica de la partícula
% Problema 09: Abrazadera en varilla horizontal con anclaje móvil
% Origen: Ayudantía 27, 2026
% =========================================================================
\subsubsection{Abrazadera en varilla horizontal con anclaje móvil}

Una abrazadera $P$ de masa $m$ desliza sin roce a lo largo de una varilla rígida, horizontal y fija (Figura \ref{fig:abrazadera}). Está unida a un resorte lineal de constante $k$ y largo natural $l_0$, cuyo otro extremo $Q$ se mueve con una ley prescrita $q(t)$. Tanto $q$ como la posición $x$ de la abrazadera se miden desde la pared izquierda, y $P$ se mantiene a la derecha de $Q$. No hay gravedad.

Resuelva usando trabajo y energía:
\begin{enumerate}
    \item Trate la fuerza del resorte como una fuerza externa (no conservativa) y obtenga la ecuación diferencial de $x(t)$ con el teorema de trabajo y energía.
    \item Trate el resorte como conservativo, con energía potencial $V(x,t)$. Calcule $dE/dt$ para $E=T+V$ y obtenga nuevamente la ecuación diferencial.
    \item Determine la potencia que debe entregar el mecanismo que mueve $Q$ y compárela con $dE/dt$.
\end{enumerate}

\begin{figure}[!ht]
    \centering
    \begin{tikzpicture}[>=Latex, scale=1.2]
        % --- PARÁMETROS ---
        \def\Lv{6.0}      % largo de la varilla [u]
        \def\xQ{1.9}      % posición del anclaje: q [u]
        \def\xP{4.5}      % posición de la abrazadera: x [u]

        % --- VARILLA Y PAREDES ---
        \draw[line width=4pt, gray!35] (0,0) -- (\Lv,0);
        \fill[pattern=north east lines] (-0.2,-0.45) rectangle (0,0.45);
        \draw[thick] (0,-0.45) -- (0,0.45);
        \draw[thick] (\Lv,-0.45) -- (\Lv,0.45);
        \fill[pattern=north east lines] (\Lv,-0.45) rectangle ({\Lv+0.2},0.45);

        % --- RESORTE, ANCLAJE Q Y ABRAZADERA P ---
        \draw[decorate, decoration={coil, aspect=0.4, segment length=3pt, amplitude=3pt}, thick] (\xQ,0) -- (\xP,0);
        \shadedraw[ball color=white!90] (\xQ,0) circle (0.13);
        \shadedraw[ball color=gray!30] (\xP,0) circle (0.13);
        \node[below=5pt] at (\xQ,0) {$Q$};
        \node[below=5pt] at (\xP,0) {$P$};
        \node[above=4pt] at ({(\xQ+\xP)/2},0) {$k$};
        \node[above=3pt, font=\small] at (\xP,0.02) {$m$};

        % --- COTAS ---
        \draw[|->, thick, gray] (0,0.75) -- (\xQ,0.75) node[midway, above] {$q(t)$};
        \draw[|-|, thick, gray] (\xQ,0.75) -- (\xP,0.75) node[midway, above] {$x-q$};
        \draw[|->, thick, gray] (0,-0.85) -- (\xP,-0.85) node[midway, below] {$x$};
        \draw[thin, gray] (\xQ,0.45) -- (\xQ,0.9);
        \draw[thin, gray] (\xP,-0.5) -- (\xP,-1.0);
    \end{tikzpicture}
    \caption{Abrazadera $P$ con anclaje móvil $Q$ y resorte lineal.}
    \label{fig:abrazadera}
\end{figure}

% --- SOLUCIÓN ---------------------------------------------------------------
\begin{solucion}{Abrazadera en varilla horizontal con anclaje móvil}

\textbf{Marcos de referencia y Sistemas de Coordenadas}\\
Basta un marco inercial fijo a la varilla, con $\vec{e}_x$ a lo largo de ella y origen en la pared izquierda.

\subsubsection*{Cinemática de la partícula}
Como $\vec{e}_x$ es constante, las posiciones se derivan directamente:
$$\vec{r}_P=x\vec{e}_x\qquad\vec{r}_Q=q\vec{e}_x\qquad\vec{v}_P=\dot{x}\vec{e}_x\qquad\vec{a}_P=\ddot{x}\vec{e}_x$$
El largo del resorte es $x-q$, su elongación es $\delta=x-q-l_0$ y su rapidez de deformación es $\dot\delta=\dot{x}-\dot{q}$. La abrazadera solo se traslada, sin rotar, así que
$$T=\frac{1}{2}m\dot{x}^2$$

\subsubsection*{Cinética}
Sobre la abrazadera actúan el resorte y la reacción de la varilla. La reacción es perpendicular a $\vec{e}_x$, o sea a la velocidad, y no trabaja: la única fuerza que trabaja es la del resorte.

\textbf{Diagrama de cuerpo libre}
\begin{center}
    \begin{tikzpicture}[>=Latex, scale=1.2]
        \draw[line width=4pt, gray!35] (-1.5,0) -- (2.5,0);
        \shadedraw[ball color=gray!30] (0,0) circle (0.14);
        \draw[->, thick, myred] (0,0) -- (-1.2,0) node[midway, above] {$k\delta$};
        \draw[->, thick, myred] (0,0) -- (0,0.9) node[above] {$N$};
        \draw[->, thick, azul] (0,-0.3) -- (1.2,-0.3) node[midway, below] {$\dot{x}$};
    \end{tikzpicture}
\end{center}
La fuerza del resorte sobre $P$ vale $\vec{F}_s=-k\delta\,\vec{e}_x$ (en la figura, con $\delta>0$, tira hacia $Q$).

\textbf{1) Resorte como fuerza no conservativa}\\
La energía cinética cambia solo por la potencia que entrega el resorte:
$$\frac{dT}{dt}=\vec{F}_s\cdot\vec{v}_P\quad\Longrightarrow\quad m\dot{x}\ddot{x}=-k\left(x-q-l_0\right)\dot{x}$$
Simplificando $\dot{x}$ (distinto de cero mientras la abrazadera se mueve):
\respuesta{$\displaystyle m\ddot{x}+k\left(x-q-l_0\right)=0$}

\textbf{2) Resorte como fuerza conservativa}\\
Incluimos el resorte en el potencial. Como el anclaje se mueve, $V$ depende explícitamente del tiempo:
$$V(x,t)=\frac{1}{2}k\left(x-q(t)-l_0\right)^2$$
Entonces $E=T+V$ no se conserva: su variación es la derivada parcial del potencial respecto al tiempo,
$$\frac{dE}{dt}=\frac{\partial V}{\partial t}=-k\left(x-q-l_0\right)\dot{q}$$
Desarrollando el lado izquierdo,
$$m\dot{x}\ddot{x}+k\left(x-q-l_0\right)\left(\dot{x}-\dot{q}\right)=-k\left(x-q-l_0\right)\dot{q}$$
los términos con $\dot{q}$ se cancelan y queda $m\dot{x}\ddot{x}+k\left(x-q-l_0\right)\dot{x}=0$. Al simplificar $\dot{x}$ resulta la misma ecuación:
\respuesta{$\displaystyle m\ddot{x}+k\left(x-q-l_0\right)=0\quad\text{es decir}\quad m\ddot{x}+kx=k\left(q(t)+l_0\right)$}
Es un oscilador de frecuencia natural $\omega_n=\sqrt{k/m}$ forzado por el anclaje. En la variable $\delta$ queda $\ddot\delta+\omega_n^2\delta=-\ddot{q}$: lo que excita al resorte es la \emph{aceleración} de $Q$.

% NOTA PARA REVISIÓN (agente B): el ítem 3 (potencia del mecanismo) NO está resuelto en Ay27S.tex;
% el desarrollo de abajo es una derivación propia del agente, consistente con el dE/dt de arriba.
\textbf{3) Potencia del mecanismo que mueve $Q$}\\
El resorte tira de $Q$ con $+k\delta\,\vec{e}_x$ (hacia $P$), de modo que el mecanismo debe ejercer $-k\delta\,\vec{e}_x$ para imponer $q(t)$. Su potencia es
$$\mathcal{P}_{mec}=\left(-k\delta\,\vec{e}_x\right)\cdot\dot{q}\,\vec{e}_x=-k\left(x-q-l_0\right)\dot{q}$$
que coincide con $dE/dt=\partial V/\partial t$: la energía de $E=T+V$ varía solo por lo que entrega (o absorbe) el mecanismo que mueve $Q$.

\end{solucion}
